\chapter{Standards}
\label{Chapter_Standards}
\lhead{Standards}

The system has been developed in accordance with the guidelines set out by India's Wireless Planning \& Coordination (WPC) Wing, specifically concerning the operation of radio frequency systems. 

The RF system is operating within the \textbf{433} MHz License-Free ISM (Industrial, Scientific, and Medical) band. 

Input Power to the RF System is limited to less than \textbf{15} dBm as per WPC guidelines, which limits the maximum allowable power level for license free transmission in this band. 

The RF System and Antenna are both designed to operate at a standard \textbf{50} $\Omega$ Input Impedance to ensure that they will be able to interoperate with other devices or equipment having the same \textbf{50} $\Omega$ input impedance. 

Performance characteristics of the Antenna have been evaluated using the standard S\textbf{11} (Return Loss) and VSWR (Voltage Standing Wave Ratio) metrics. 

Bluetooth Low Energy (BLE) wireless protocol is used to facilitate Node-to-Phone wireless communication. 

LoRa (Long Range) Modulation Layer is used as the Physical Modulation Layer to support low data rate, long range wireless communication. The use of LoRa as the Physical Modulation Layer supports the implementation of long range, low power consumption wireless networking. 

CSMA/CA (Carrier Sense Multiple Access/Collision Avoidance) Protocol is used to manage channel access. 

Network Logic and Packet Structure are based on an open source protocol called Meshtastic. 

AES-GCM (Advanced Encryption Standard in Galois/Counter Mode) Protocol, recognized by NIST as a Standard for Authenticated Encryption, is used to secure all communications.
